%%% ==========================================================================
%%%  第四章：紧性
%%%  提纲与两轮审阅意见见 temp/ch04-outline.md，读书笔记见
%%%  math/notes/ch04-compactness.md。写作时须守住的几条（均为审阅指出）：
%%%   - 闭单位球不紧的反例用 g_n = x^(2^n)，不可用 x^n（后者下标相邻时
%%%     距离趋于零，不是 ε-分离的）；
%%%   - 序列紧属拓扑层，不可与全有界并列在度量层；
%%%   - Lebesgue 数引理须由序列紧证出，不得借道最值定理，否则循环论证；
%%%   - 最值定理须写「非空」紧集；
%%%   - 「闭且有界即紧」只对欧氏空间成立，「有限维」的说法只在赋范空间中才准确。
%%%  TODO：第 7、9 章与卷四、卷五建立后，把「第 N 章」改回 \cref{ch:...}。
%%% ==========================================================================
\chapter{紧性}
\label{ch:compactness}

\begin{keypoint}[本章要点]
  \begin{itemize}
    \item 连续性本身推不出「取到最值」。缺的那条性质是定义域的紧性，
          它把无限的覆盖化为有限的覆盖。
    \item 紧性同样有三副面孔：序列紧、全有界加完备、开覆盖紧。
          与\cref{ch:continuity} 的连续性一样，只有不提距离的那一副
          活到卷四，并在那里反客为主。
    \item 在任何度量空间中，紧集必闭且有界；逆命题在欧氏空间中成立，
          在一般度量空间中不成立（两个反例，成因各异）。
    \item 第 1 章说紧性是确界原理在拓扑层面的接替者。本章把这句话坐实：
          \([a,b]\) 紧要用确界原理来证，此后凡涉及一般度量空间的论证都改用紧性。
          （实值函数的最值仍要用到 \(\R\) 的确界性质，见\cref{thm:extreme-value}。）
  \end{itemize}
\end{keypoint}

\section{为什么需要紧性}
\label{sec:why-compact}

第 2 章定义 \(C[a,b]\) 上的一致度量时写的是
\begin{equation}
  \label{eq:uniform-metric-recall}
  d_{\infty}(f,g) = \max_{a \le x \le b} \abs{f(x)-g(x)} ,
\end{equation}
当时附了一句：最大值确实存在，这一点将在本章由紧性给出，此处暂且承认
（\cref{ex:metric-CX}）。现在来还这笔账。

先看清欠的究竟是什么。取 \(f(x)=x\)，它在 \((0,1)\) 上连续，
上确界为 \(1\) 却取不到；取 \(f(x)=\arctan x\)，它在 \(\R\) 上连续，
上确界为 \(\pi/2\) 也取不到。两个函数都无可指摘，问题出在定义域：
一个少了端点，一个伸向无穷。

可见「连续函数取到最值」不是连续性的推论，而依赖定义域的某种性质。
本章要找出那条性质，它叫紧性。

\begin{concept}{紧性}
  \cwhy{连续性本身推不出取到最值。上面两个例子表明，缺的是定义域的一条性质；
    本章要把它找出来并证明它足够。}
  \crel{它是第 1 章确界原理在拓扑层面的接替者。第 1 章的确界原理依赖序，
    离开 \(\R\) 就无从谈起；紧性只用开集就能说清，因而走得更远。}
  \cwhere{分析学与拓扑学的枢纽。凡是要把「局部成立」升格为「整体成立」的地方
    都用它。}
  \cdo{把无限化为有限：从局部控制所对应的开覆盖中选出\emph{有限}子覆盖，
    再对有限多个正数取最小值。最值定理、一致连续定理、
    以及第 7 章 Riemann 可积性的判别，用的都是这一步。}
\end{concept}

\begin{pitfall}
  「化无限为有限」是对开覆盖那一副面孔的说法，不可滥用。
  无穷多个正数\emph{未必}有最小值，其下确界还可能为 \(0\)；
  必须先取出有限子覆盖，才谈得上取最小。同样，
  「无穷数列取收敛子列」也并没有把无限变成有限。
\end{pitfall}

紧性也有三副面孔，这与\cref{ch:continuity} 的连续性同属一个结构。

\begin{center}
\begin{tabular}{@{}lll@{}}
  \toprule
  面孔 & 长处 & 到卷四 \\
  \midrule
  序列紧（\cref{sec:seq-compact}） & 承接第 3 章，证明最省力 & 仍有定义，与紧不再等价 \\
  全有界加完备（\cref{sec:totally-bounded}） & 最便于检验，可具体估算 & 依赖度量，无从谈起 \\
  开覆盖紧（\cref{sec:open-cover}） & 只用开集 & 取作定义，与序列紧不再等价 \\
  \bottomrule
\end{tabular}
\end{center}

两章的平行不是巧合。但要把话说准：到一般拓扑空间中，开覆盖紧与序列紧
\emph{都}仍有定义（收敛是拓扑概念，见\cref{prop:convergence-topological}），
变的是两者不再普遍等价；全有界与完备性则需要度量，无从谈起。
这是全书「用得越少，走得越远」这条主线的第二次出现。

\subsection*{本章要还的账}

前几章有三处写着「第 4 章会解释」，但它们的性质并不相同。
只有第一处真正靠紧性结清，另两处是用来划定定理边界的。

\begin{center}
\begin{tabular}{@{}llp{4.4cm}@{}}
  \toprule
  出处 & 结论 & 靠什么 \\
  \midrule
  \cref{ex:metric-CX} & \(C[a,b]\) 中 \(\max\) 存在 & \textbf{紧性}，\cref{thm:extreme-value} 结清 \\
  \cref{ex:not-uniform} & \(1/x\) 在 \((0,1]\) 上不一致连续 & 两列反例，并非「不紧」所致 \\
  \cref{ex:one-over-x-uniform} & \(1/x\) 在 \([a,\infty)\)（\(a>0\)）上一致连续 & Lipschitz 估计；该区间并不紧 \\
  \bottomrule
\end{tabular}
\end{center}

\begin{remark}
  \label{rem:sufficient-not-necessary}
  这张表要先看明白，否则容易把\cref{thm:uniform-on-compact} 记反。
  该定理说的是：\emph{定义域紧}是「连续蕴含一致连续」的\emph{充分}条件。
  它不说定义域不紧就没有一致连续的映射，\cref{ex:one-over-x-uniform}
  正是反例；也不说定义域不紧就必有连续而不一致连续的映射，
  尽管\cref{ex:not-uniform} 确实找到了一个。
\end{remark}

\begin{table}[htbp]
  \centering
  \caption{本章各条结论在后续章节中的用途。}
  \label{tab:ch4-where-used}
  \begin{tabular}{@{}lll@{}}
    \toprule
    本章结论 & 首次使用 & 用途 \\
    \midrule
    最值定理       & 第 6 章 & 中值定理族的前提 \\
    一致连续定理   & 第 7 章 & Riemann 可积性的判别 \\
    连续像仍紧     & 第 5 章 & 连通性的对应结论 \\
    全有界         & 第 9 章 & Arzelà--Ascoli 定理 \\
    开覆盖紧       & 卷四     & 取作拓扑空间中紧性的定义 \\
    \bottomrule
  \end{tabular}
\end{table}

\begin{remark}[读法建议]
  \cref{sec:seq-compact} 与\cref{sec:consequences} 必须掌握。
  \cref{sec:open-cover} 的三者等价性证明初读可略，记住结论即可；
  其中的 Lebesgue 数引理值得单独看一遍，它的手法在后文反复出现。
  \cref{sec:totally-bounded} 的全有界在本卷用于第 9 章的 Arzelà--Ascoli 定理
  （见\cref{tab:ch4-where-used}），到卷五则是通往紧算子理论的门。
\end{remark}

\section{序列紧}
\label{sec:seq-compact}

依\cref{ch:continuity} 的先例，先用序列作定义，因为序列是此刻手上最熟的工具。

\begin{definition}[序列紧]
  \label{def:seq-compact}
  设 \((X,d)\) 为度量空间，\(A \subseteq X\)。若 \(A\) 中每个数列都有
  收敛于 \(A\) 中某点的子列，则称 \(A\) 为\term{序列紧的}{sequentially compact}。
\end{definition}

\begin{example}[平凡的情形]
  \label{ex:trivial-compact}
  空集与有限集序列紧。
  \begin{worked}
    空集中没有数列，条件空洞地成立。

    设 \(A\) 有限，\(\seq{a_{n}}\) 为 \(A\) 中的数列。由抽屉原理，
    \(A\) 中至少有一个元素 \(p\) 被取到无穷多次，取这些下标构成子列，
    则该子列恒为 \(p\)，收敛于 \(p \in A\)。
  \end{worked}
\end{example}

\begin{theorem}
  \label{thm:compact-closed-bounded}
  序列紧集必闭且有界。
\end{theorem}

\begin{proof}
  \emph{闭}。设 \(p \in \overline{A}\)。由\cref{exr:sequential-closure}，
  存在 \(A\) 中数列 \(a_{n} \to p\)。由序列紧，某子列 \(a_{n_{k}} \to q \in A\)。
  而子列的下标满足 \(n_{k} \ge k\)，故 \(a_{n} \to p\) 蕴含 \(a_{n_{k}} \to p\)；
  极限唯一给出 \(q = p\)，于是 \(p \in A\)。故 \(A = \overline{A}\)。

  \emph{有界}。固定 \(p \in X\)。若 \(A\) 无界，则对每个 \(n\) 存在
  \(a_{n} \in A\) 使 \(d(p,a_{n}) \ge n\)。由序列紧取收敛子列 \(a_{n_{k}} \to q\)，
  则 \(\seq{a_{n_{k}}}\) 收敛故为 Cauchy 列、故有界
  （\cref{prop:cauchy-basics} 之 (i)(ii)），
  与 \(d(p,a_{n_{k}}) \ge n_{k} \to \infty\) 矛盾。
\end{proof}

逆命题一般不成立。这是本章最该防的误解，两个反例的成因不同。

\begin{example}[闭且有界而不紧，其一：不完备]
  \label{ex:q-closed-bounded}
  在度量空间 \(\Q\)（取 \(\R\) 的子空间度量）中，
  \(A = \setst{x \in \Q}{0 \le x \le 2,\ x^{2} \le 2}\) 闭且有界，但不序列紧。
  \begin{worked}
    \(A\) 在 \(\Q\) 中闭：若 \(q \in \Q\) 且 \(q \notin A\)，则 \(q<0\) 或 \(q^{2}>2\)，
    两种情形都能取到与 \(A\) 不交的球。有界显然。

    取\cref{eq:sqrt2-decimal} 的数列 \(1, 1.4, 1.41, \dots\)，各项都在 \(A\) 中。
    它在 \(\R\) 中收敛于 \(\sqrt{2}\)，故其任一子列也收敛于 \(\sqrt{2}\)；
    而 \(\sqrt{2} \notin \Q\)，故没有子列收敛于 \(A\) 中的点。

    毛病出在 \(\Q\) 不完备。
  \end{worked}
\end{example}

\begin{example}[闭且有界而不紧，其二：不全有界]
  \label{ex:ball-not-compact}
  \(C[0,1]\)（取一致度量）中的闭单位球
  \(B = \setst{f}{d_{\infty}(f,0) \le 1}\) 闭且有界，但不序列紧。
  \begin{worked}
    取 \(g_{n}(x) = x^{2^{n}}\)。各 \(g_{n} \in B\)。设 \(m > n\)，
    取 \(x_{0}\) 使 \(x_{0}^{2^{n}} = 1/2\)，则
    \(x_{0}^{2^{m}} = \bigl(x_{0}^{2^{n}}\bigr)^{2^{m-n}} = (1/2)^{2^{m-n}} \le 1/4\)，
    于是
    \begin{equation}
      \label{eq:separated}
      d_{\infty}(g_{n},g_{m}) \ge \abs{g_{n}(x_{0}) - g_{m}(x_{0})}
        \ge \frac{1}{2} - \frac{1}{4} = \frac{1}{4}.
    \end{equation}
    诸 \(g_{n}\) 两两相距至少 \(1/4\)，故任一子列都不是 Cauchy 列，
    更不收敛。故 \(B\) 不序列紧。
  \end{worked}
\end{example}

\begin{pitfall}
  \begin{enumerate}[label=(\arabic*), leftmargin=1.6em]
    \item \textbf{以为「闭且有界」就是紧。} 上面两个反例说明不然。
          这条等价只在欧氏空间中成立（\cref{thm:heine-borel}）。
    \item \textbf{误称 \(f_{n}(x)=x^{n}\) 的各项一致分离。}
          \(\norm{x-x^{2}}_{\infty} = 1/4\)，而 \(\norm{x^{100}-x^{101}}_{\infty} \approx 0.0037\)，
          下标相邻时距离趋于零，故整列并非一致分离。
          \(f_{n}\) 本身\emph{可以}用来证闭单位球不序列紧——任一子列的逐点极限
          在 \([0,1)\) 上为 \(0\)、在 \(1\) 处为 \(1\)，不连续，故不一致收敛。
          要证\emph{不全有界}则须找出一致分离的无限子族，
          \cref{ex:ball-not-compact} 所取的 \(x^{2^{n}}\) 正是 \(x^{n}\) 的一个这样的子列。
    \item \textbf{把\cref{ex:ball-not-compact} 说成「有限维特有」而不加限定。}
          准确的陈述是：在\emph{赋范线性空间}中，闭单位球紧当且仅当空间有限维
          （卷五的 Riesz 引理）。一般度量空间没有维数可言。
  \end{enumerate}
\end{pitfall}

\section{全有界与完备}
\label{sec:totally-bounded}

\cref{ex:q-closed-bounded} 与\cref{ex:ball-not-compact} 各自缺了一条性质：
前者缺完备性，后者缺一种「用有限多个小球盖得住」的性质。本节说明
这两样合起来恰好等于序列紧。

\begin{definition}[全有界]
  \label{def:totally-bounded}
  设 \((X,d)\) 为度量空间。若对每个 \(\varepsilon>0\)，存在\emph{有限个}点
  \(x_{1},\dots,x_{N} \in X\) 使
  \begin{equation}
    \label{eq:eps-net}
    X = \bigcup_{j=1}^{N} B(x_{j},\varepsilon),
  \end{equation}
  则称 \(X\) 为\term{全有界的}{totally bounded}，
  称 \(\set{x_{1},\dots,x_{N}}\) 为一个 \term{\(\varepsilon\)-网}{epsilon-net}。
  这里的「网」指有限的逼近点集，与\cref{ch:continuity} 中以有向集为指标的
  拓扑网是两个不同的概念，只是中译名相同。
\end{definition}

\begin{intuition}
  有界是「装得进一个大球」，全有界是「用有限多个\emph{任意小}的球盖得住」。
  后者严格更强。在 \(\R^{n}\) 中两者一致，因为一个大方块总能切成有限多个小方块；
  在无穷维空间中则不然：\cref{ex:ball-not-compact} 中诸 \(g_{n}\) 两两相距 \(1/4\)，
  半径 \(1/8\) 的球每个至多装一个，而它们有无穷多个。
  （这说的是\emph{整个}闭单位球；无穷维空间中当然也有全有界的子集，
  例如任一有限集。）
\end{intuition}

\cref{def:totally-bounded} 要求网点落在空间自身之内。
把它用到子空间 \(A \subseteq X\) 上时，网点就得落在 \(A\) 里，
而现成的覆盖往往是用 \(X\) 中的点作球心造出来的。下面这条引理补上这一步，
代价只是把半径放大一倍。

\begin{lemma}[网点可移入子集]
  \label{lem:net-into-subset}
  设 \(A \subseteq X\)。若对每个 \(\varepsilon>0\) 都存在有限多个
  \(x_{1},\dots,x_{N} \in X\) 使 \(A \subseteq \bigcup_{j} B(x_{j},\varepsilon)\)，
  则 \(A\) 作为度量子空间全有界。
\end{lemma}

\begin{proof}
  给定 \(\varepsilon>0\)，先取球心在 \(X\) 中的 \(\varepsilon/2\)-覆盖
  \(A \subseteq \bigcup_{j=1}^{N} B(x_{j},\varepsilon/2)\)。
  丢掉那些与 \(A\) 不交的球；对余下的每个 \(j\) 取一点
  \(a_{j} \in A \cap B(x_{j},\varepsilon/2)\)。

  任给 \(a \in A\)，它落在某个保留下来的 \(B(x_{j},\varepsilon/2)\) 中，于是
  \[
    d(a,a_{j}) \le d(a,x_{j}) + d(x_{j},a_{j}) < \varepsilon/2 + \varepsilon/2 = \varepsilon .
  \]
  故诸 \(a_{j}\) 构成 \(A\) 中的有限 \(\varepsilon\)-网。
\end{proof}

\begin{figure}[htbp]
  \centering
  \begin{tikzpicture}[scale=0.86]
    \def\blob{plot[smooth cycle, tension=0.7] coordinates
      {(-1.0,0.1) (-0.4,0.9) (0.25,0.5) (0.95,0.8) (0.8,-0.3) (0.1,-0.4) (-0.35,-0.9)}}
    \foreach \dx/\ttl in {0/{有界}, 4.9/{全有界}, 9.8/{有界而不全有界}} {
      \node[below] at (\dx,-2.35) {\small \ttl};
    }
    % (甲) 装得进一个大球
    \draw[htGray, dashed] (0,0) circle (1.9);
    \draw[htRule, very thick] \blob;
    \draw[htGray, -{Stealth[length=4pt]}] (0,0) -- (135:1.9);
    \node[htGray, font=\scriptsize, above right] at (135:1.1) {\(R\)};
    \fill[htGray] (0,0) circle (0.04);
    % (乙) 有限多个小球盖得住
    \begin{scope}[xshift=4.9cm]
      \foreach \cx/\cy in {-0.78/0.05, -0.3/0.62, 0.2/0.42, 0.72/0.62,
                           0.72/-0.12, 0.24/-0.22, -0.22/-0.66, -0.1/0.1} {
        \fill[htGray, opacity=0.18] (\cx,\cy) circle (0.52);
        \draw[htGray, thin] (\cx,\cy) circle (0.52);
      }
      \draw[htRule, very thick] \blob;
      \draw[htGray, -{Stealth[length=4pt]}] (0.72,-0.12) -- ++(-35:0.52);
      \fill[htGray] (0.72,-0.12) circle (0.035);
      \node[htGray, font=\scriptsize, below right=-1pt and -2pt] at (0.95,-0.3) {\(\varepsilon\)};
    \end{scope}
    % (丙) 两两分离，有限多个小球盖不住
    \begin{scope}[xshift=9.8cm]
      \draw[htGray, dashed] (0,0) circle (1.9);
      \foreach \ang in {40, 90, 140, 190, 240, 290} {
        \draw[htGray, thin] (\ang:1.25) circle (0.34);
        \fill[htRule] (\ang:1.25) circle (0.055);
      }
      \node[font=\footnotesize, htGray] at (335:1.25) {\(\cdots\)};
      \draw[htRule, {Stealth[length=4pt]}-{Stealth[length=4pt]}] (90:1.25) -- (140:1.25);
      \node[htRule, font=\scriptsize] at (115:0.72) {\(\ge \tfrac{1}{4}\)};
    \end{scope}
  \end{tikzpicture}
  \caption{有界与全有界。左：整个集合装得进一个半径 \(R\) 的球。
    中：对每个 \(\varepsilon>0\)，都能用\emph{有限}多个半径 \(\varepsilon\) 的球盖住。
    右：\cref{ex:ball-not-compact} 的情形——集合有界，但其中有无穷多个点两两相距
    至少 \(1/4\)，半径 \(1/8\) 的球每个至多装一个，于是有限多个盖不住。
    在 \(\R^{n}\) 中左、中两栏等价（\cref{exr:bounded-vs-tb}），
    在一般度量空间中不等价。}
  \label{fig:bounded-vs-tb}
\end{figure}

\begin{proposition}
  \label{prop:tb-implies-bounded}
  全有界蕴含有界，反之不然。
\end{proposition}

\begin{proof}
  取 \(\varepsilon=1\) 的有限网 \(x_{1},\dots,x_{N}\)，
  令 \(R = 1 + \max_{j} d(x_{1},x_{j})\)，则 \(X \subseteq B(x_{1},R)\)。

  反例：无限集配离散度量，有界（全体含于半径 \(2\) 的球内）
  而不全有界（\(\varepsilon = 1/2\) 时每个球只含一点，需无穷多个）。
\end{proof}

\begin{theorem}
  \label{thm:seq-compact-iff}
  度量空间 \((X,d)\) 序列紧，当且仅当它完备且全有界。
\end{theorem}

\begin{strategy}
  必要性的两半分开做：完备性由「Cauchy 列有收敛子列则收敛」立得；
  全有界则用反证——若某个 \(\varepsilon\) 没有有限网，就能逐点挑出
  两两相距不小于 \(\varepsilon\) 的一列点，它没有收敛子列。

  充分性要造子列。手法是\emph{逐层加细}：先用 \(1\)-网把数列分到有限多个球里，
  必有一个球含无穷多项，取出这一段作第一层子列；再用 \(1/2\)-网对这一段
  重做一次，得第二层；如此下去。最后取\emph{对角线}——第 \(j\) 层的第 \(j\) 项。
  这个对角线子列的项彼此越挨越近，故为 Cauchy 列，由完备性收敛。
\end{strategy}

\begin{proof}
  \emph{必要性}。设 \(X\) 序列紧。若 \(\seq{a_{n}}\) 是 Cauchy 列，
  由序列紧它有收敛子列，再由\cref{prop:cauchy-basics} 之 (iii)
  知它收敛，故 \(X\) 完备。

  设 \(X\) 不全有界，则存在 \(\varepsilon_{0}>0\) 使 \(X\) 没有有限 \(\varepsilon_{0}\)-网。
  任取 \(x_{1} \in X\)。因 \(B(x_{1},\varepsilon_{0}) \neq X\)，可取 \(x_{2}\) 使
  \(d(x_{1},x_{2}) \ge \varepsilon_{0}\)。归纳地，已取 \(x_{1},\dots,x_{n}\) 之后，
  因这 \(n\) 个球不覆盖 \(X\)，可取 \(x_{n+1} \notin \bigcup_{j \le n} B(x_{j},\varepsilon_{0})\)。
  所得数列满足 \(d(x_{m},x_{n}) \ge \varepsilon_{0}\)（\(m \neq n\)），
  任一子列都不是 Cauchy 列，故不收敛。这与序列紧矛盾。

  \emph{充分性}。设 \(X\) 完备且全有界，\(\seq{a_{n}}\) 为 \(X\) 中的数列。

  取 \(1\)-网。有限多个球覆盖 \(X\)，故其中至少一个含有 \(\seq{a_{n}}\) 的无穷多项；
  把落在该球中的那些项按原顺序排出，记作 \(\seq{a^{(1)}_{n}}\)，它是 \(\seq{a_{n}}\)
  的子列，且任意两项相距小于 \(2\)。

  对 \(\seq{a^{(1)}_{n}}\) 用 \(1/2\)-网重做一次，得子列 \(\seq{a^{(2)}_{n}}\)，
  任意两项相距小于 \(1\)。如此下去，得一列层层嵌套的子列 \(\seq{a^{(j)}_{n}}\)，
  第 \(j\) 层中任意两项相距小于 \(2/j\)。

  取\emph{对角线} \(b_{j} = a^{(j)}_{j}\)。因各层嵌套，\(b_{j}\) 在原数列中的下标
  随 \(j\) 严格递增，故 \(\seq{b_{j}}\) 确是 \(\seq{a_{n}}\) 的子列。
  又当 \(k \ge j\) 时 \(b_{k}\) 属于第 \(j\) 层，故 \(d(b_{j},b_{k}) < 2/j\)，
  于是 \(\seq{b_{j}}\) 是 Cauchy 列。由完备性它收敛。
\end{proof}

\begin{remark}
  \label{rem:diagonal-new}
  上面的\emph{对角线抽取}是本书第一次使用这个手法，值得记住：
  当需要一个同时满足无穷多个条件的子列时，就逐层加细、最后取对角线。
  它在第 9 章证明 Arzelà--Ascoli 定理时还会出现。
\end{remark}

\section{开覆盖紧}
\label{sec:open-cover}

第三副面孔的陈述里没有距离。

\begin{definition}[开覆盖与紧]
  \label{def:compact}
  设 \((X,d)\) 为度量空间，\(A \subseteq X\)。称一族开集 \(\set{U_{i}}_{i \in I}\)
  为 \(A\) 的\term{开覆盖}{open cover}，若 \(A \subseteq \bigcup_{i \in I} U_{i}\)。
  若 \(A\) 的每个开覆盖都有\emph{有限}子覆盖，即存在有限的
  \(i_{1},\dots,i_{N} \in I\) 使 \(A \subseteq \bigcup_{k=1}^{N} U_{i_{k}}\)，
  则称 \(A\) 为\term{紧的}{compact}。
\end{definition}

\begin{example}
  \label{ex:not-compact-cover}
  \((0,1)\) 不紧。
  \begin{worked}
    取开覆盖 \(U_{n} = (1/n,\ 1)\)，\(n \ge 2\)。它确实覆盖 \((0,1)\)：
    每个 \(x \in (0,1)\) 都有 \(n\) 使 \(1/n < x\)。
    但任取有限多个 \(U_{n_{1}},\dots,U_{n_{k}}\)，
    令 \(N = \max_{j} n_{j}\)，则它们的并为 \((1/N,1)\)，
    漏掉了 \(1/(2N)\)。故没有有限子覆盖。

    对比\cref{ex:trivial-compact}：紧与不紧的差别，
    在这个例子里表现为「能否用有限多块盖住」。
  \end{worked}
\end{example}

下面这条引理是本节的技术核心，也是「从序列控制转向有限覆盖」的关口。

\begin{figure}[htbp]
  \centering
  \begin{tikzpicture}[x=8.2cm, y=1cm]
    % 上：有限 ε-网
    \draw[htGray, thick] (0,2.9) -- (1,2.9);
    \foreach \j/\lb in {0/{0}, 1/{\frac{1}{5}}, 2/{\frac{2}{5}}, 3/{\frac{3}{5}}, 4/{\frac{4}{5}}, 5/{1}} {
      \pgfmathsetmacro\xj{\j/5}
      \pgfmathsetmacro\yj{3.24 + 0.28*mod(\j,2)}
      \fill[htRule] (\xj,2.9) circle (0.6pt);
      \draw[htGray, thin, dotted] (\xj,2.9) -- (\xj,\yj);
      \draw[htRule, {|[width=3.2pt]}-{|[width=3.2pt]}] (\xj-0.1,\yj) -- (\xj+0.1,\yj);
      \node[below, font=\scriptsize, htGray] at (\xj,2.84) {\(\lb\)};
    }
    \node[left, font=\scriptsize] at (-0.012,2.9) {\([0,1]\)};
    \node[right, font=\scriptsize, align=left] at (1.015,3.5)
      {\(N+1=6\) 个网点\\ 每段半宽 \(\varepsilon>\tfrac{1}{10}\)};

    % 下：Lebesgue 数
    \fill[htGray, opacity=0.15] (0.4,0.7) rectangle (0.6,1.66);
    \draw[htRule, very thick] (0,1.54) -- (0.6,1.54);
    \node[right, font=\scriptsize] at (0.615,1.54) {\(U_{1}=[0,\,0.6)\)};
    \draw[htRule, very thick] (0.4,1.18) -- (1,1.18);
    \node[right, font=\scriptsize] at (1.015,1.18) {\(U_{2}=(0.4,\,1]\)};
    \draw[htGray, thick] (0,0.82) -- (1,0.82);
    \foreach \x/\lb in {0/{0}, 0.4/{0.4}, 0.5/{0.5}, 0.6/{0.6}, 1/{1}} {
      \draw[htGray] (\x,0.76) -- (\x,0.88);
      \node[below, font=\scriptsize, htGray] at (\x,0.74) {\(\lb\)};
    }
    \draw[htRule, {|[width=3.2pt]}-{|[width=3.2pt]}, very thick] (0.4,0.2) -- (0.6,0.2);
    \fill[htRule] (0.5,0.2) circle (0.6pt);
    \node[below, font=\scriptsize] at (0.5,0.13) {\(B(0.5,\,\delta)\)，\(\delta=0.1\)};
    \node[right, font=\scriptsize, align=left] at (1.015,0.2)
      {重叠区 \((0.4,0.6)\)\\ 定出最大的 \(\delta\)};
  \end{tikzpicture}
  \caption{上：\([0,1]\) 的有限 \(\varepsilon\)-网（\cref{exr:eps-net-count}），
    \(N+1\) 个点的 \(\varepsilon\)-球盖住整个区间。
    下：覆盖 \(\set{U_{1},U_{2}}\) 的 Lebesgue 数（\cref{lem:lebesgue}）。
    最坏的点是 \(x=0.5\)：往右要留在 \(U_{1}\) 内需 \(\delta \le 0.1\)，
    往左要留在 \(U_{2}\) 内也需 \(\delta \le 0.1\)。可用的 \(\delta\)
    由两片开集的\emph{重叠}宽度定，而非由单片的长度定。
    这也是\cref{lem:lebesgue} 的证明为何要在极限点处取球的原因。}
  \label{fig:net-and-lebesgue}
\end{figure}


\begin{lemma}[Lebesgue 数]
  \label{lem:lebesgue}
  设 \((X,d)\) 序列紧，\(\set{U_{i}}_{i \in I}\) 为 \(X\) 的开覆盖。
  则存在 \(\delta>0\)，使每个半径为 \(\delta\) 的球都含于某一个 \(U_{i}\) 之中。
  这样的 \(\delta\) 称为该覆盖的 \term{Lebesgue 数}{Lebesgue number}。
\end{lemma}

\begin{proof}
  反设不存在。则对每个 \(n\)，\(\delta = 1/n\) 不合用，
  即存在 \(x_{n} \in X\) 使 \(B(x_{n},1/n)\) 不含于任何一个 \(U_{i}\)。

  由序列紧，取子列 \(x_{n_{k}} \to x\)。因诸 \(U_{i}\) 覆盖 \(X\)，
  存在 \(i_{0}\) 使 \(x \in U_{i_{0}}\)；又 \(U_{i_{0}}\) 开，
  存在 \(r>0\) 使 \(B(x,r) \subseteq U_{i_{0}}\)。

  取 \(k\) 充分大，使 \(d(x_{n_{k}},x) < r/2\) 且 \(1/n_{k} < r/2\)。
  则对任意 \(y \in B(x_{n_{k}},1/n_{k})\)，
  \begin{equation}
    \label{eq:lebesgue-est}
    d(y,x) \le d(y,x_{n_{k}}) + d(x_{n_{k}},x) < \frac{r}{2} + \frac{r}{2} = r ,
  \end{equation}
  故 \(B(x_{n_{k}},1/n_{k}) \subseteq B(x,r) \subseteq U_{i_{0}}\)，
  与 \(x_{n_{k}}\) 的取法矛盾。
\end{proof}

\begin{pitfall}
  \cref{lem:lebesgue} 只能由序列紧证出，\emph{不可}借道
  「紧集上的连续函数取到正的最小值」。后者是\cref{thm:extreme-value}
  的推论，而\cref{thm:extreme-value} 又要用本节的\cref{thm:compact-equiv}，
  绕回去就成了循环论证。本书中凡遇到「甲由乙证、乙由甲证」的嫌疑，
  都应把依赖顺序显式写出来查一遍。
\end{pitfall}

\begin{theorem}[三者等价]
  \label{thm:compact-equiv}
  设 \((X,d)\) 为度量空间。下述三条等价：
  \begin{enumerate}[label=(\roman*)]
    \item \(X\) 紧（每个开覆盖有有限子覆盖）；
    \item \(X\) 序列紧；
    \item \(X\) 完备且全有界。
  \end{enumerate}
\end{theorem}

\begin{proof}
  (ii) 与 (iii) 的等价即\cref{thm:seq-compact-iff}。以下证
  (i) \(\Rightarrow\) (ii) \(\Rightarrow\) (i)。

  \emph{(i) \(\Rightarrow\) (ii)}。设 \(X\) 紧而 \(\seq{a_{n}}\) 没有收敛子列。
  此处要用的是\emph{数列}的聚点：称 \(x\) 为 \(\seq{a_{n}}\) 的聚点，
  若每个 \(B(x,r)\) 都含有无穷多个下标 \(n\) 所对应的项 \(a_{n}\)。
  这与\cref{ch:metric} 中\emph{集合}的聚点（用去心球定义）不是一回事：
  常值数列 \(a_{n}=x\) 以 \(x\) 为数列的聚点，而单点集 \(\set{x}\) 没有集合聚点。
  \(x\) 是数列聚点当且仅当有子列收敛于 \(x\)（取 \(r=1/k\) 逐次抽取即可）。

  于是由反设，每个 \(x \in X\) 都不是 \(\seq{a_{n}}\) 的聚点，
  故存在 \(r_{x}>0\) 使 \(B(x,r_{x})\) 只含 \(\seq{a_{n}}\) 的有限多项。
  诸 \(B(x,r_{x})\) 构成开覆盖，由紧性取有限子覆盖
  \(B(x_{1},r_{x_{1}}),\dots,B(x_{N},r_{x_{N}})\)。
  它们各含有限多项而并为 \(X\)，故数列总共只有有限多项，矛盾。

  \emph{(ii) \(\Rightarrow\) (i)}。设 \(X\) 序列紧，\(\set{U_{i}}\) 为开覆盖。
  由\cref{lem:lebesgue} 取 Lebesgue 数 \(\delta>0\)。
  由\cref{thm:seq-compact-iff}，\(X\) 全有界，故存在有限的 \(\delta\)-网
  \(x_{1},\dots,x_{N}\)，即 \(X = \bigcup_{j} B(x_{j},\delta)\)。
  每个 \(B(x_{j},\delta)\) 含于某个 \(U_{i_{j}}\)，
  于是 \(X = \bigcup_{j=1}^{N} U_{i_{j}}\)，得有限子覆盖。
\end{proof}

\begin{remark}[通往卷四的门]
  \label{rem:compact-to-topology}
  \cref{def:compact} 的陈述中只出现「开集」，没有距离。
  卷四把它取作拓扑空间中紧性的定义。届时\cref{thm:compact-equiv} 不再成立：
  一般拓扑空间中紧与序列紧互不蕴含，全有界更是无从谈起（那里没有度量）。
  这与\cref{ch:continuity} 的情形一致：三副面孔中，提到距离的那一副
  （全有界）离开度量空间即无从谈起，而序列紧虽仍有定义，
  却不再与紧等价。只有开覆盖紧原封不动地活了下来并反客为主。

  文献中还有第四种刻画，称为可数紧：每个\emph{可数}开覆盖有有限子覆盖。
  在度量空间中它与上述三条等价，本讲义不用，故不展开。
\end{remark}

\section{Heine--Borel 定理}
\label{sec:heine-borel}

现在回到 \(\R\) 与 \(\R^{n}\)。这里「闭且有界」与紧确实等价，
而证明的第一步要用到第 1 章的确界原理。

\begin{theorem}
  \label{thm:interval-compact}
  闭区间 \([a,b] \subseteq \R\) 序列紧。
\end{theorem}

\begin{proof}
  设 \(\seq{x_{n}}\) 为 \([a,b]\) 中的数列。令
  \begin{equation}
    \label{eq:pugh-set}
    C = \setst{x \in [a,b]}{x_{n} < x \text{ 只对有限多个 } n \text{ 成立}} .
  \end{equation}
  \(a \in C\)（无 \(x_{n} < a\)），故 \(C\) 非空；\(b\) 是 \(C\) 的上界。
  由第 1 章的确界原理，\(c = \sup C\) 存在且 \(c \in [a,b]\)。

  往证对每个 \(\varepsilon>0\)，\((c-\varepsilon,\,c+\varepsilon)\) 中含有 \(\seq{x_{n}}\) 的无穷多项。

  \emph{下侧}。由 \(c=\sup C\)，取 \(c' \in C\) 满足 \(c'>c-\varepsilon\)。
  按 \(C\) 的定义，\(x_{n}<c'\) 只有限多次，故 \(x_{n} \le c-\varepsilon\) 也只有限多次，
  即最终 \(x_{n} > c-\varepsilon\)。

  \emph{上侧}分两种情形。若 \(c+\varepsilon>b\)，则一切 \(x_{n} \le b < c+\varepsilon\)。
  若 \(c+\varepsilon \le b\)，则 \(c+\varepsilon \in [a,b]\) 而 \(c+\varepsilon>c=\sup C\)，
  故 \(c+\varepsilon \notin C\)，按定义即有无穷多个 \(n\) 使 \(x_{n}<c+\varepsilon\)。

  两侧合起来：\((c-\varepsilon,\,c+\varepsilon)\) 中有无穷多项。
  依次取 \(\varepsilon = 1, 1/2, 1/3, \dots\) 并保证下标递增，即得收敛于 \(c\) 的子列。
  这个写法同时覆盖了 \(c=a\)、\(c=b\) 与退化区间的情形。
\end{proof}

\begin{remark}
  \label{rem:sup-to-compact}
  这个证明兑现了第 1 章那句话：紧性是确界原理的接替者。
  交接的确切含义是：\emph{定义域}的那一侧从此改用紧性——
  \cref{thm:uniform-on-compact} 对任意紧度量空间成立，不必是实区间。
  但\emph{值域}若为 \(\R\)，其确界性质仍要用到：
  \cref{thm:extreme-value} 的证明里就构造了 \(\sup f(X)\)。
  紧性接管的是定义域，不是实数系本身。
\end{remark}

\begin{theorem}[Heine--Borel]
  \label{thm:heine-borel}
  \(A \subseteq \R^{n}\) 紧，当且仅当 \(A\) 闭且有界。
\end{theorem}

\begin{proof}
  必要性即\cref{thm:compact-closed-bounded}。

  充分性。设 \(A\) 闭且有界，\(\seq{a^{(m)}}\) 为 \(A\) 中的数列，
  记 \(a^{(m)} = \bigl(a^{(m)}_{1},\dots,a^{(m)}_{n}\bigr)\)。
  由有界性，存在 \(M\) 使各坐标满足 \(\abs{a^{(m)}_{j}} \le M\)。

  逐坐标抽取。第一坐标构成 \([-M,M]\) 中的数列，
  由\cref{thm:interval-compact} 有收敛子列；在该子列上再看第二坐标，
  又有收敛子列；如此进行 \(n\) 次。所得的最终子列使\emph{每个}坐标都收敛，
  由第 2 章关于 \(\R^{n}\) 中收敛的刻画（逐坐标收敛即收敛），
  该子列在 \(\R^{n}\) 中收敛于某点 \(p\)。

  因 \(A\) 闭而子列各项在 \(A\) 中，故 \(p \in A\)。于是 \(A\) 序列紧，
  由\cref{thm:compact-equiv} 即紧。
\end{proof}

\begin{remark}
要点不在「只能抽有限次」——可数多个坐标也可以用对角线抽取同时处理
  （\cref{thm:seq-compact-iff} 的证明就是这么做的）。要点在于：
  有限维时，\emph{有限}多个坐标的收敛足以控制范数；无穷维时不然。
  \(\ell^{2}\) 中的标准基 \(e_{n}\) 逐坐标趋于 \(0\)，而 \(\norm{e_{n}}_{2}=1\) 恒不变。
  请回看\cref{ex:q-closed-bounded} 与\cref{ex:ball-not-compact}：
  本定理对欧氏空间成立，对一般度量空间不成立。
\end{remark}

\section{紧性的三条推论}
\label{sec:consequences}

第 1 章曾说，闭区间上连续函数的几条经典结论应当统一归因于紧性。
本节兑现这句话。三条结论的证明都很短，因为力气已经花在\cref{thm:compact-equiv} 上了。

\begin{theorem}[连续像仍紧]
  \label{thm:continuous-image-compact}
  设 \(f \mapsdef X \to Y\) 连续，\(X\) 紧，则 \(f(X)\) 紧。
\end{theorem}

\begin{proof}
  用序列判据。设 \(\seq{y_{n}}\) 为 \(f(X)\) 中的数列，取 \(x_{n} \in X\) 使
  \(f(x_{n}) = y_{n}\)。由 \(X\) 序列紧，某子列 \(x_{n_{k}} \to x \in X\)。
  由 \(f\) 连续（\cref{def:continuous}），\(y_{n_{k}} = f(x_{n_{k}}) \to f(x) \in f(X)\)。
\end{proof}

\begin{theorem}[最值定理]
  \label{thm:extreme-value}
  设 \(X\) 为\emph{非空}紧度量空间，\(f \mapsdef X \to \R\) 连续。
  则 \(f\) 取到最大值与最小值。
\end{theorem}

\begin{proof}
  由\cref{thm:continuous-image-compact}，\(f(X) \subseteq \R\) 紧，
  由\cref{thm:compact-closed-bounded} 它闭且有界。
  因 \(X\) 非空，\(f(X)\) 非空，故 \(\beta = \sup f(X)\) 存在（第 1 章确界原理）。
  由上确界的 \(\varepsilon\) 刻画，存在 \(f(X)\) 中数列趋于 \(\beta\)，
  故 \(\beta \in \overline{f(X)} = f(X)\)，即存在 \(x^{*} \in X\) 使 \(f(x^{*}) = \beta\)。
  最小值同理。
\end{proof}

\begin{remark}
  「非空」不可省：空集紧（\cref{ex:trivial-compact}），
  而空集上的函数无最值可言。
\end{remark}

这就结清了第一笔账：\(C[a,b]\) 中的 \(d_{\infty}\) 之所以能写成 \(\max\)，
是因为 \([a,b]\) 紧（\cref{thm:interval-compact}）而 \(\abs{f-g}\) 连续，
由\cref{thm:extreme-value} 最大值确实取到。

\begin{theorem}[一致连续定理]
  \label{thm:uniform-on-compact}
  设 \(X\) 紧，\(f \mapsdef X \to Y\) 连续，则 \(f\) 一致连续。
\end{theorem}

\begin{strategy}
  用\cref{thm:uniform-seq} 的两列判据反证最省力：若不一致连续，
  就有两列点彼此越挨越近而像不接近；紧性给出收敛子列，
  连续性再逼出矛盾。这条路只用序列，不必碰 \(\varepsilon\)-\(\delta\) 的两层量词。
\end{strategy}

\begin{proof}
  设 \(f\) 不一致连续。由\cref{thm:uniform-seq}，
  存在两列 \(\seq{x_{n}}\)、\(\seq{y_{n}}\) 与 \(\varepsilon_{0}>0\)，使
  \begin{equation}
    \label{eq:uc-contradiction}
    d(x_{n},y_{n}) \to 0, \qquad
    \rho\bigl(f(x_{n}),f(y_{n})\bigr) \ge \varepsilon_{0} .
  \end{equation}
  由 \(X\) 序列紧，取子列 \(x_{n_{k}} \to x\)。
  由 \(d(x_{n_{k}},y_{n_{k}}) \to 0\) 得 \(y_{n_{k}} \to x\) 亦成立。
  于是由 \(f\) 连续，\(f(x_{n_{k}}) \to f(x)\) 且 \(f(y_{n_{k}}) \to f(x)\)，
  故 \(\rho(f(x_{n_{k}}),f(y_{n_{k}})) \to 0\)，与\cref{eq:uc-contradiction} 矛盾。
\end{proof}

\begin{remark}[边界，务必看清]
  \label{rem:uc-boundary}
  \cref{thm:uniform-on-compact} 说的是：\emph{定义域紧}是「连续蕴含一致连续」的
  \emph{充分}条件。它不说别的。对照\cref{sec:why-compact} 的那张表：
  \begin{itemize}
    \item \cref{ex:one-over-x-uniform} 中 \([a,\infty)\)（\(a>0\)）\emph{不紧}
          （无界），而 \(1/x\) 在其上仍一致连续——靠的是 Lipschitz 估计
          \(\abs{1/x-1/y} \le \abs{x-y}/a^{2}\)，不是本定理。
          故定义域紧并非一致连续的必要条件。
    \item \cref{ex:not-uniform} 中 \((0,1]\) 不紧，而 \(1/x\) 在其上确实
          不一致连续——但这是由两列点算出来的，\emph{不能}仅凭「不紧」推出。
          本定理只有一个方向。
  \end{itemize}
  这两条都不是本定理的推论；列在此处是为了划出它的边界。
\end{remark}

另有一条推论留作习题：紧集上的连续双射必为同胚（\cref{exr:compact-homeo}）。
一般而言连续双射未必同胚——由\cref{exr:discrete-continuous}，
从 \(([0,1],\,\text{离散度量})\) 到 \(([0,1],\,\text{通常度量})\) 的恒等映射
连续且是双射，而其逆不连续。差别正在于定义域紧与否：离散度量下的
\([0,1]\) 不紧（其中收敛列必最终为常数，而 \(\seq{1/n}\) 无收敛子列）。

\section{本章结论在更一般框架下的地位}
\label{sec:ch4-outlook}

\begin{table}[htbp]
  \centering
  \caption{紧性相关概念所依赖的结构。须区分两件事：一个概念\emph{依赖}什么结构，
    与两个概念在何处\emph{等价}。序列紧只涉及序列收敛，而收敛由拓扑决定
    （\cref{prop:convergence-topological}），故它属拓扑层；但它与开覆盖紧的等价
    只在度量空间中成立。}
  \label{tab:ch4-structure}
  \begin{tabular}{@{}lll@{}}
    \toprule
    层次 & 概念 & 一般拓扑空间中的情况 \\
    \midrule
    拓扑 & 开覆盖紧、序列紧 & 两者一般\emph{不}等价 \\
    度量 & 全有界、完备性   & 不能仅由开集族决定 \\
    序   & 上确界性质       & 需另行指定序结构 \\
    \bottomrule
  \end{tabular}
\end{table}

三条向前的接口：

\emph{其一}，\cref{def:compact} 在卷四取作拓扑空间中紧性的定义。
届时会有紧而不序列紧的空间，也有序列紧而不紧的空间；
把\cref{thm:compact-equiv} 拆开来看，正是那里的主要工作之一。
Tychonoff 定理断言任意多个紧空间之积仍紧，本讲义在卷四证明。

\emph{其二}，\cref{def:totally-bounded} 的全有界在第 9 章与卷五重现。
Arzelà--Ascoli 定理以「有界加等度连续」刻画 \(C(K)\) 中的\emph{相对紧}子集，
即闭包紧者；要断言集合自身紧，还须加上闭性。
（反例：\(K=[0,1]\) 上全体常值函数 \(f_{c} \equiv c\)，\(0<c<1\)。
该族一致有界且等度连续，但 \(f_{1/n}\) 趋于不属于该族的零函数，故不闭、不紧。）
其证明的核心是本章的对角线抽取（见\cref{thm:seq-compact-iff} 之后的注）。

\emph{其三}，\cref{thm:uniform-on-compact} 在第 7 章用来证明
闭区间上的连续函数 Riemann 可积。那里需要把区间分割得足够细，
使函数在每个小段上的振幅一致地小；这正是一致连续所保证的。

\section*{习题}
\addcontentsline{toc}{section}{习题}

\begin{exercise}[例行]
  \label{exr:discrete-compact}
  在离散度量空间中，哪些子集紧？
  \begin{solution}
    恰是有限子集。有限集紧见\cref{ex:trivial-compact}。
    反之若 \(A\) 无限，取其中两两不同的一列点，
    则任意两项距离为 \(1\)，无收敛子列，故不紧。

    对照\cref{prop:tb-implies-bounded}：无限离散空间有界而不全有界，
    由\cref{thm:seq-compact-iff} 即知不紧。
  \end{solution}
\end{exercise}

\begin{exercise}[例行]
  \label{exr:closed-subset-compact}
  证明紧集的闭子集紧；两个紧集之并紧。
  \begin{solution}
    设 \(A\) 紧、\(F \subseteq A\) 闭，\(\seq{a_{n}}\) 为 \(F\) 中的数列。
    由 \(A\) 序列紧取子列 \(a_{n_{k}} \to p \in A\)；因 \(F\) 闭且各项在 \(F\) 中，
    \(p \in F\)。故 \(F\) 序列紧。

    设 \(A, B\) 紧，\(\seq{c_{n}}\) 为 \(A \cup B\) 中的数列。
    \(A\) 与 \(B\) 至少有一个含 \(\seq{c_{n}}\) 的无穷多项，
    对那一段用序列紧即得。
  \end{solution}
\end{exercise}

\begin{exercise}[例行]
  \label{exr:eps-net-count}
  为 \([0,1]\) 构造 \(\varepsilon\)-网并估算所需点数；再对 \([0,1]^{n}\) 做同样的事。
  \begin{solution}
    取整数 \(N > 1/(2\varepsilon)\)，令 \(x_{j} = j/N\)（\(j=0,\dots,N\)）。
    任一 \(x \in [0,1]\) 与最近的 \(x_{j}\) 相距不超过 \(1/(2N) < \varepsilon\)，
    故这 \(N+1\) 个点构成 \(\varepsilon\)-网。

    \([0,1]^{n}\)（取 \(d_{\infty}(x,y)=\max_{i}\abs{x_{i}-y_{i}}\)）：
    取各坐标网格点的乘积 \(\setst{(j_{1}/N,\dots,j_{n}/N)}{0 \le j_{i} \le N}\)，
    共 \((N+1)^{n}\) 个点。任一 \(x\) 的每个坐标与最近的 \(j_{i}/N\) 相距不超过
    \(1/(2N)\)，取最大值仍不超过 \(1/(2N)<\varepsilon\)，故这是 \(\varepsilon\)-网。

    点数随维数指数增长，这个现象在高维数值计算中称为维数灾难。
    但它说明的只是网点数会爆炸：每个\emph{有限}维立方体仍然全有界。
    \cref{ex:ball-not-compact} 中失效的是无穷维空间的整个闭单位球，
    并非那里的每个有界子集都不全有界。
  \end{solution}
\end{exercise}

\begin{exercise}[例行]
  \label{exr:lebesgue-number-concrete}
  设 \(X = [0,1]\)，开覆盖取 \(U_{1} = [0,\,0.6)\)、\(U_{2} = (0.4,\,1]\)。
  求一个可用的 Lebesgue 数。
  \begin{solution}
    取 \(\delta = 0.1\)。若 \(x < 0.5\)，则 \(B(x,\delta) \cap X \subseteq [0,\,x+0.1)
    \subseteq [0,0.6) = U_{1}\)；若 \(x \ge 0.5\)，则 \(B(x,\delta) \cap X
    \subseteq (x-0.1,\,1] \subseteq (0.4,1] = U_{2}\)。故 \(\delta=0.1\) 可用。

    它还是最大的可用值：在 \(x=0.5\) 处，\(B(0.5,\delta) \subseteq U_{1}\) 要求
    \(\delta \le 0.1\)，\(B(0.5,\delta) \subseteq U_{2}\) 也要求 \(\delta \le 0.1\)。
    可用的 \(\delta\) 由两片开集的\emph{重叠区间} \((0.4,0.6)\) 定，
    而非由单个成员的长度定。
  \end{solution}
\end{exercise}

\begin{exercise}[例行]
  \label{exr:arctan-no-max}
  证明 \(\arctan\) 在 \(\R\) 上连续、有界，却取不到上确界。
  这与\cref{thm:extreme-value} 矛盾吗？
  \begin{solution}
    \(\arctan\) 是 \(\R\) 到 \((-\pi/2,\,\pi/2)\) 上的双射，故连续且
    \(\abs{\arctan x} < \pi/2\) 恒成立，即 \(\pi/2\) 是上界而取不到。
    它还是\emph{最小}的上界：任给 \(t<\pi/2\)，因 \(\arctan\) 满射到
    \((-\pi/2,\pi/2)\)，存在 \(x\) 使 \(\arctan x \in (t,\,\pi/2)\)，
    故 \(t\) 不是上界。于是 \(\sup_{\R}\arctan = \pi/2\)，而它不在值域内。

    不矛盾：\(\R\) 不紧（无界，由\cref{thm:compact-closed-bounded}）。
    \cref{thm:extreme-value} 的前提不满足。
  \end{solution}
\end{exercise}

\begin{exercise}[结构]
  \label{exr:fip}
  证明：\(X\) 紧当且仅当任一具有\emph{有限交性质}的闭集族
  （即其中任意有限多个成员之交非空）其总交非空。
  \begin{solution}
    取补集即可。设 \(X\) 紧，\(\mathcal{F}\) 为闭集族且总交为空。
    则 \(\setst{X \setminus F}{F \in \mathcal{F}}\) 是开覆盖，
    取有限子覆盖 \(X = \bigcup_{j}(X \setminus F_{j}) = X \setminus \bigcap_{j}F_{j}\)，
    故 \(\bigcap_{j} F_{j} = \emptyset\)，与有限交性质矛盾。
    反向同理。

    这个形式在卷四证明 Tychonoff 定理时比开覆盖形式好用。
  \end{solution}
\end{exercise}

\begin{exercise}[结构]
  \label{exr:bounded-vs-tb}
  比较有界与全有界。给出一个有界而不全有界的度量空间，
  并说明在 \(\R^{n}\) 中两者为何一致。
  \begin{solution}
    无限集配离散度量：全体含于半径 \(2\) 的球内故有界；
    而 \(\varepsilon=1/2\) 的球只含一点，需无穷多个，故不全有界。

    设 \(A \subseteq \R^{n}\) 有界，则 \(A \subseteq [-M,M]^{n}\)。
    由\cref{exr:eps-net-count}（把 \([0,1]^{n}\) 伸缩平移到 \([-M,M]^{n}\)），
    该立方体有有限 \(\varepsilon\)-网，其点未必在 \(A\) 中；
    由\cref{lem:net-into-subset}，\(A\) 自身全有界。
    关键在于\emph{有限}维：坐标只有 \(n\) 个，网格点数才是有限的。
  \end{solution}
\end{exercise}

\begin{exercise}[结构]
  \label{exr:lebesgue-alt}
  用 Lebesgue 数引理另证\cref{thm:uniform-on-compact}，
  并比较两种证法各自用到了什么。
  \begin{solution}
    给定 \(\varepsilon>0\)。由 \(f\) 连续，每个 \(p \in X\) 有 \(\delta_{p}>0\) 使
    \(f\bigl(B(p,\delta_{p})\bigr) \subseteq B\bigl(f(p),\varepsilon/2\bigr)\)。
    诸 \(B(p,\delta_{p}/2)\) 构成开覆盖，取 Lebesgue 数 \(\delta\)。
    若 \(d(x,y)<\delta\)，则 \(x,y\) 同属某个 \(B(p,\delta_{p}/2)\)，
    故 \(\rho(f(x),f(y)) \le \rho(f(x),f(p)) + \rho(f(p),f(y)) < \varepsilon\)。

    比较：正文的证法只用序列判据与序列紧，不必碰两层量词；
    本证法用开覆盖紧与 Lebesgue 数，但给出了 \(\delta\) 的来源，
    更接近可计算的形式。两条路对应本章的两副面孔。
  \end{solution}
\end{exercise}

\begin{exercise}[结构]
  \label{exr:compact-separable}
  证明紧度量空间可分，即含有可数的稠密子集。
  \begin{solution}
    对每个 \(n\) 取有限的 \(1/n\)-网 \(G_{n}\)（由\cref{thm:seq-compact-iff}
    全有界）。令 \(G = \bigcup_{n} G_{n}\)，它是可数多个有限集之并，故可数。
    任给 \(x\) 与 \(\varepsilon>0\)，取 \(n > 1/\varepsilon\)，
    则 \(G_{n}\) 中有点与 \(x\) 相距小于 \(1/n < \varepsilon\)。故 \(G\) 稠密。
  \end{solution}
\end{exercise}

\begin{exercise}[结构]
  \label{exr:why-not-general}
  \cref{thm:heine-borel} 对 \(\R^{n}\) 成立而对一般度量空间不成立。
  指出证明中哪一步无法推广。
  \begin{solution}
    「逐坐标抽取子列」那一步。它要求坐标只有\emph{有限}个，
    才能在有限次抽取之后同时控制住全部坐标。
    \(C[0,1]\) 中没有有限多个坐标可抽，
    \cref{ex:ball-not-compact} 正是这一失效的具体表现。

    另一处是\cref{ex:q-closed-bounded}：\(\Q\) 中闭且有界推不出完备，
    而\cref{thm:seq-compact-iff} 要求完备。
  \end{solution}
\end{exercise}

\begin{exercise}[延伸]
  \label{exr:compact-homeo}
  设 \(X\) 紧，\(f \mapsdef X \to Y\) 为连续双射，\(Y\) 为度量空间。
  证明 \(f\) 是同胚。再与\cref{sec:consequences} 末尾的反例对照。
  \begin{solution}
    只需证 \(f^{-1}\) 连续，等价于 \(f\) 把闭集映成闭集
    （\(f\) 双射时 \((f^{-1})^{-1}(F) = f(F)\)）。
    设 \(F \subseteq X\) 闭，由\cref{exr:closed-subset-compact} 它紧；
    由\cref{thm:continuous-image-compact}，\(f(F)\) 紧；
    由\cref{thm:compact-closed-bounded}，\(f(F)\) 闭。

    对照：\([0,1]\) 从离散度量到通常度量的恒等映射是连续双射而非同胚，
    那里的定义域不紧。紧性正是补上这个缺口的条件。
    （在一般拓扑空间中，本题须要求陪域为 Hausdorff 空间。）
  \end{solution}
\end{exercise}

\begin{exercise}[延伸]
  \label{exr:riesz-preview}
  设 \(E\) 为赋范空间，\(\seq{e_{n}}\) 为其中一列满足
  \(\norm{e_{n}}=1\) 且 \(\norm{e_{m}-e_{n}} \ge 1/2\)（\(m \neq n\)）的元素。
  证明 \(E\) 的闭单位球不紧。再说明这与\cref{ex:ball-not-compact} 的关系。
  \begin{solution}
    诸 \(e_{n}\) 在闭单位球内且两两相距不小于 \(1/2\)，
    故任一子列都不是 Cauchy 列，闭单位球不序列紧。

    \cref{ex:ball-not-compact} 中的 \(g_{n}(x)=x^{2^{n}}\) 正是这样一列
    （相距 \(\ge 1/4\)，把 \(1/2\) 换成 \(1/4\) 即可）。
    卷五的 Riesz 引理断言：在赋范空间中这样的一列存在，
    当且仅当空间无限维。于是「闭单位球紧」恰好刻画了有限维。
  \end{solution}
\end{exercise}

\begin{exercise}[延伸]
  \label{exr:cantor-compact}
  Cantor 集 \(K\) 由 \([0,1]\) 反复挖去中间三分之一开区间所得。
  证明 \(K\) 紧。
  \begin{solution}
    记 \(K_{0}=[0,1]\)，\(K_{n+1}\) 为从 \(K_{n}\) 的每个闭区间挖去
    中间三分之一开区间所得，\(K = \bigcap_{n} K_{n}\)。

    每个 \(K_{n}\) 是有限多个闭区间之并，故闭；可数多个闭集之交仍闭，
    故 \(K\) 闭。又 \(K \subseteq [0,1]\) 有界。
    由\cref{thm:heine-borel}，\(K\) 紧。

    \(K\) 将在第 7 章重现：它不可数却是零测集，
    这个反差是 Riemann 积分与 Lebesgue 积分之别的第一个信号。
  \end{solution}
\end{exercise}
